% !TeX root = ../复分析笔记.tex
% 由主文档通过 \input 载入；请编译主文档。
\setcounter{chapter}{2} % 保留原书章号。
\chapter{综合例题与方法选择}
\label{chap:ca-synthesis}

\section{同一道积分的三种理解}

\begin{example}[路径半径改变时发生了什么]
对 $r>0$、$r\ne1$，计算正向圆积分
\[
 I(r)=\oint_{|z|=r}\frac{\mathrm dz}{z^2(1-z)}.
\]
\end{example}
\begin{solve}
先列出奇点：原点为二阶极点，$1$ 为简单极点。对 $0<r<1$，在去心单位圆盘内展开
\[
 \frac1{z^2(1-z)}=z^{-2}+z^{-1}+1+z+\cdots,
\]
所以原点留数为 $1$，积分 $I(r)=2\pi i$。同一结果也可直接由高阶柯西公式得到：取 $f(z)=1/(1-z)$，则 $\oint f(z)/z^2\,\mathrm dz=2\pi i f'(0)=2\pi i$。

当 $r>1$ 时，还要计入 $z=1$ 的留数 $-1$，因此 $I(r)=2\pi i(1-1)=0$。若改从大圆环展开，
\[
 \frac1{z^2(1-z)}=-\sum_{n=0}^\infty z^{-n-3},\qquad |z|>1,
\]
其中没有 $z^{-1}$ 项，也给出零。三种方法相互印证，并说明积分在路径连续变形且不跨越奇点时保持不变；半径穿过 $1$ 时恰好越过一个奇点。$r=1$ 的路径经过极点，普通曲线积分不存在，不能通过左右结果插值得到。
\end{solve}

\section{把增长、唯一性与局部数据结合起来}

\begin{example}[少量数据何时决定整个函数]
设 $f$ 为整函数，存在常数 $C>0$ 使 $|f(z)|\le C(1+|z|)$，且 $f(0)=0$、$f(1)=1$。求 $f$。
\end{example}
\begin{solve}
增长条件配合柯西估计使 $f^{(n)}(0)=0$ 对所有 $n\ge2$ 成立，所以 $f(z)=az+b$。两个点值给出 $b=0$、$a=1$，故 $f(z)=z$。两个点值本身不能决定任意整函数，例如 $z+cz(z-1)$ 对任意 $c\in\mathbb C$ 都满足这两个点值；真正把自由度压缩到一次多项式的是整体增长条件。
\end{solve}

\section{选择方法前应核对的条件}

遇到复积分，先标出定义域、奇点与路径方向，再决定能否用原函数、柯西公式或留数定理；遇到级数，先确定中心与圆盘或圆环，再选择等比展开的方向；遇到映射，先给出分支及值域，再证明单射；遇到边值问题，先明确边界数据与解的函数类别，再谈存在性和唯一性。下表把常用结论与最容易遗漏的条件放在一起，供复习时查阅。

\begin{longtable}{@{}p{0.20\textwidth}p{0.35\textwidth}p{0.37\textwidth}@{}}
\toprule
工具 & 必须核对的条件 & 结论与使用边界 \\
\midrule
\endhead
C--R 方程 & 点处实可微，或用连续偏导作为充分保证 & 得到点处复可导；只有方程成立不够。 \\
原函数端点公式 & 同一个定义域上存在单值原函数 & 路径积分仅依赖端点；$1/z$ 须检查对数分支。 \\
柯西积分定理 & 函数在相关路径与所需内部全纯，或使用单连通版本 & 闭路积分为零；路径上的全纯性不能代替内部条件。 \\
柯西积分公式 & 分子全纯，求值点在内部，曲线正向且不经过该点 & 分母 $n+1$ 次幂对应 $n$ 阶导数与因子 $1/n!$。 \\
泰勒展开 & 中心附近全纯，确定可延拓的圆盘 & 只有收敛范围内才能逐项微分、积分并表示原函数。 \\
洛朗展开 & 在指定圆环上单值全纯 & 系数依赖圆环；分类中心奇点须使用完整去心圆盘。 \\
留数定理 & 路径不经过奇点，区域内孤立奇点满足定理范围 & 正向简单闭曲线取留数和；一般曲线还要考虑绕数。 \\
鲁歇定理 & 两函数在边界及内部附近全纯，边界严格模不等式 & 保持零点总数，按重数计；不直接给出各根位置。 \\
恒等定理 & 两函数在同一区域全纯，相等点在内部有聚点 & 可推出全域相同；聚点只在边界时不能这样使用。 \\
最大模原理 & 开集上的全纯性及相应连通性 & 非常数函数的模不能在内部局部取最大；最小模需无零点。 \\
共形映射 & 局部要求导数非零，整体还需单射并确定值域 & 保角不等于保距，局部可逆不等于整体可逆。 \\
泊松公式 & 连续实值圆周数据；内部求调和解 & 给出连续边值问题唯一解，不要求任意复值数据都可全纯延拓。 \\
\bottomrule
\end{longtable}

这些方法背后有一条统一主线：复可导性限制局部变化，柯西理论把局部性质传递到积分，级数把积分信息转成系数，系数再决定零点、奇点与映射结构。真正掌握课程，不只是记住最后的计算公式，还要能说明一个条件失效时，证明在哪一步中断，并能给出相应反例。


